\section{Conclusion}\label{sec:conclusion}

When a graph recommender is evaluated on items inserted uniformly at random into feeds, two log-only signals help it rank those items. Lifelong, exposure-conditional affinities of users for content keys (LEAF) transfer from items outside the test pool through content keys only, and the gain disappears when the users' lifelong tables are shuffled. A detection-adjusted item attractiveness (DAX), fitted by EM to a click model that discounts non-clicks by user and scenario examination, adds a further increment on top of EASE, the QP family, DIN and LEAF-LGN. The increment vanishes when attractiveness is shuffled across items, while within-tag shuffles keep part of it.

On KuaiRand-27K, under a temporal protocol with five user-level splits, LEAF-LGN with DAX is higher than each of the seven comparators of the main family that have results, with Holm-adjusted per-user t tests below $10^{-11}$ and split-level Nadeau-Bengio Holm-adjusted p values of at most $4.56\times10^{-2}$. The comparison with the QP family with DAX, which is the fair baseline in this study, gives $+12.04$ [1.63]. The result holds on KuaiRand-1K and across three training seeds for LEAF-LGN with DAX.

The evidence is limited to one platform family, five splits, and search budgets that differ by two orders of magnitude. The claims are the controlled ranking gain under the stated protocol and the module's increment on several backbones. The paper does not claim unbiasedness, causal identification of item quality, or superiority over lifelong deep models tuned more widely, and the SIM comparison remains to be run.
